\subsection{Bag of Words (BoW)}

Le modèle \textit{Bag of Words} représente chaque tweet comme un vecteur de comptage sur un
vocabulaire $\mathcal{V}$ de 20\,000 termes. Cette approche ignore l'ordre des mots mais
constitue une baseline rapide et interprétable.

\subsection{TF-IDF}

Le score TF-IDF pondère chaque terme $t$ dans un document $d$ selon :
\begin{equation}
  \text{TF-IDF}(t,d) = \log(1 + \text{tf}(t,d)) \times \log\frac{N}{\text{df}(t)}
\end{equation}
Nous utilisons \texttt{sublinear\_tf=True} et 20\,000 features maximum.

\subsection{N-grams}

En étendant le TF-IDF aux bigrammes et trigrammes (\texttt{ngram\_range=(1,3)}),
nous capturons des expressions multi-mots porteuses de sentiment comme \textit{not good}
ou \textit{very happy}. Le vocabulaire passe à 30\,000 features.

\subsection{Word2Vec}

Word2Vec (Skip-Gram, dimension 200) entraîne des vecteurs denses sur notre corpus.
Chaque tweet est représenté par la \textbf{moyenne} des vecteurs de ses tokens.
Cette approche capture des relations sémantiques : $\text{vec}(\textit{good}) \approx \text{vec}(\textit{great})$.

\subsection{BERT Embeddings}

DistilBERT génère des représentations \textbf{contextuelles} : le même mot a des vecteurs
différents selon son contexte. Nous fine-tunons le modèle complet pour la classification binaire.

\begin{figure}[H]
\centering
\includegraphics[width=\textwidth]{figures/03_representation.png}
\caption{Comparaison des méthodes de représentation : taille du vocabulaire, dimension
des vecteurs et densité des matrices résultantes.}
\label{fig:repr}
\end{figure}
